% Candidate extension: implemented in the standalone EGO scene-adapter branch.
% NOT part of the scene:false joint128/replay method above.
% Integration into the joint model and its evaluation are pending.
\subsection{Observed Geometry Adaptation: Candidate Extension}
\label{sec:scene_candidate}

\paragraph{Scope of this draft.}
This subsection describes the implemented standalone EGO scene adapter. Its integration with the joint EGO/EXO model remains pending; it is not a description of the current joint128 experiment.

Camera reconstruction can additionally provide observed scene points in the motion coordinate system. We formulate scene conditioning through body-relative measurements around a motion proposal. A frozen EGO denoiser first predicts a kinematically reconciled motion $X^b$. Skinning the predicted body yields 67 surface anchors $q_{t,k}$, their body-region labels, and associated global rotations $Q_{t,k}$.

For each anchor, we query nearby measured points and form a 14-dimensional relation descriptor
\begin{equation}
\begin{aligned}
r_{t,k}=\big[&
Q_{t,k}^{\top}(p^\star-q_{t,k})/r_q,\
d/r_q,\
Q_{t,k}^{\top}n,\
\eta/r_q,\\
&v_n,\ c_n,\ \rho_s,\ f,\ h,\ u
\big],
\end{aligned}
\label{eq:scene_relations}
\end{equation}
where $p^\star$ is the nearest observed point, $d$ its distance, $n$ a locally fitted normal, and $\eta=(q_{t,k}-p^\star)^\top n$ an oriented plane offset. The remaining entries encode normal validity, neighborhood coverage, surface reliability, measured free-space evidence, a nearby-point indicator, and an unknown-space indicator. The query radius is $r_q=0.6\,\mathrm{m}$. Invalid geometric entries are masked.

Free-space evidence is restricted to measured camera-to-first-hit ray segments. A local plane offset does not establish the inside/outside sign of a full scene surface. The unknown indicator marks anchors with neither a nearby observation nor measured free-space evidence; missing geometry is not converted to a free-space label.

An MLP embeds each relation together with anchor motion features, an anchor identity embedding, and the diffusion signal level. Evidence-weighted tokens are pooled within 22 body regions and processed by temporal convolutions. Zero-initialized projections inject the resulting features into the bottleneck and two decoder resolutions of a second pass through the same frozen denoiser. Each injection is further gated by
\begin{equation}
G_t=\left(\max_k g_{t,k}\right)\bar\alpha_\tau^{1/2},
\label{eq:scene_gate}
\end{equation}
where $g_{t,k}$ combines local surface reliability, a weaker point-only confidence, and measured free-space evidence.

Let $Y^b$ and $Y^g$ denote the raw baseline and geometry-conditioned outputs. The output correction is bounded before kinematic reconciliation:
\begin{equation}
\begin{aligned}
\widetilde Y^p &=Y^{b,p}
+\kappa^{-1}\mathcal T_{c_s}
\big(\kappa(Y^{g,p}-Y^{b,p})\big),\\
\widetilde Y^R&=Y^{b,R}
+c_6\tanh\big((Y^{g,R}-Y^{b,R})/c_6\big),
\end{aligned}
\label{eq:scene_bounds}
\end{equation}
where $\mathcal T_c(v)=v\,\tanh(\|v\|_2/c)/(\|v\|_2/c)$ with its continuous value at zero, $c_s=0.15\,\mathrm{m}$, and $c_6=0.5$. The rotation limit is imposed on 6D coordinates rather than rotation angles. The final positions are reconstructed through FK with the predicted rest skeleton.

This adapter currently uses a frozen EGO prior. Its proposal/query pass has a different role from the differentiable reference pass in SRRD. A joint scene model must specify their composition, gradient paths, and inference cost before these components can be presented as one evaluated method.

